% chapter1.tex - Introducción

\chapter{Orígenes de la teoría de conjuntos}

El origen de la teoría de conjuntos se remonta al último tercio del siglo XIX y se debe al matemático alemán Georg
Cantor (San Petersburgo, 1845 -- Halle, 1918). La teoría de un problema concreto del análisis: en 1870 Eduard Heine le
propuso a Cantor demostrar la unicidad de la representación de una función mediante una serie trigonométrica. Al
analizar los conjuntos excepcionales de puntos donde esa unicidad podía fallar, Cantor estudió los conjuntos de puntos
y de sus \emph{conjuntos derivados}. De ahí surgió, su artículo de 1874 \emph{Über eine Eigenschaft des Inbegriffes
    aller reellen algebraischen Zahlen}, en donde se presentó por primera vez la demostración de que los números reales no
pueden ponerse en correspondencia biunívoca con los naturales. Ese trabajo suele considerarse el acta de nacimiento de
la teoría de conjuntos. Cantor introdujo después la noción de \emph{potencia} o cardinalidad (1878), la cual permite
comparar el tamaño de conjuntos infinitos, y con ella demostró que existen distintos tamaños de infinito; en 1891 dio
de este hecho su demostración más célebre, el argumento diagonal.

Las ideas de Cantor hicieron de la noción de conjunto la principal herramienta para la fundamentación de las
matemáticas, pues permitieron formalizar conceptos como los de número, función y relación. De acuerdo con Cantor, un
conjunto es una colección, considerada como un todo, de objetos distintos y bien determinados de nuestra intuición o de
nuestro pensamiento, llamados \emph{elementos} del conjunto. En particular, los conjuntos pueden ser elementos de otros
conjuntos. Decir que un conjunto $x$ es un elemento de un conjunto $y$ se denota por el símbolo $x\in y$. Si un
conjunto $x$ no es un elemento de un conjunto $y$, se denota por el símbolo $x\notin y$.

Un conjunto que consta de un número finito de elementos $x_{1}, x_{2}, \ldots, x_{n}$ se denota por
\[
    \{x_{1}, x_{2}, \ldots, x_{n}\}.
\]
Dos conjuntos $x$ e $y$ son \textbf{iguales}\index{igualdad de conjuntos}, lo que se escribe $x = y$, si y solo si
tienen los mismos elementos; es decir, si todo elemento de $x$ es un elemento de $y$ y, recíprocamente, todo elemento
de $y$ es un elemento de $x$. Por ejemplo, los conjuntos $\{x, y\}$ y $\{y, x\}$ son iguales, ya que ambos contienen
los mismos elementos $x$ e $y$.

Un conjunto que no contiene ningún elemento se llama \textbf{conjunto vacío}\index{conjunto vacío} y se denota por
$\emptyset$. Como dos conjuntos con los mismos elementos son iguales, el conjunto vacío es único: no puede haber dos
conjuntos vacíos distintos.

\section{Los números naturales como conjuntos}

La teoría desarrollada hasta aquí permite dar un significado preciso a los números naturales, que son abstracciones
creadas por los seres humanos para facilitar el conteo. La idea, debida a John von Neumann, consiste en identificar
cada número natural con el conjunto de los números que lo preceden:

\[
    \begin{aligned}
        0 & = \emptyset,         \\
        1 & = \{0\},             \\
        2 & = \{0, 1\},          \\
        3 & = \{0, 1, 2\},       \\
        4 & = \{0, 1, 2, 3\},    \\
        5 & = \{0, 1, 2, 3, 4\}, \\
          & \ \ \vdots
    \end{aligned}
\]

En general, el número natural $n$ se identifica con el conjunto $\{0, 1, \ldots, n-1\}$ de sus predecesores, de modo
que $n$ tiene exactamente $n$ elementos, y el paso al número siguiente se obtiene mediante la operación
\textbf{sucesor}\index{sucesor}
\[
    n + 1 = n \cup \{n\}.
\]

El conjunto $\{0, 1, 2, 3, 4, 5, \ldots\}$ de todos los números naturales se denota por $\omega$. El conjunto $\{1, 2,
    3, 4, 5, \ldots\}$ de los números naturales no nulos se denota por $\mathbb{N}$.

\begin{remark}[¿Es el cero un número natural?]
    No hay un acuerdo universal sobre si el cero debe contarse entre los números naturales: se trata de una convención
    que varía de una tradición matemática a otra. Buena parte de la tradición europea continental incluye el cero,
    mientras que la tradición de Europa del Este suele empezar en el uno. Una huella cotidiana de esta discrepancia
    aparece en la numeración de los pisos de los edificios, que en unos países comienza en la planta baja y en otros
    en el primer piso. Para evitar toda ambigüedad, en este libro reservamos $\omega$ para $\{0, 1, 2, \ldots\}$ y
    $\mathbb{N}$ para $\{1, 2, 3, \ldots\}$.
\end{remark}

\section{El constructor de conjuntos}

Con mucha frecuencia necesitamos formar un conjunto reuniendo los objetos que poseen cierta propiedad; por ejemplo, el
conjunto de los números impares. Para ello disponemos del \textbf{constructor de conjuntos}\index{constructor de
    conjuntos} $\{x : \varphi(x)\}$, que denota exactamente lo que necesitamos: el conjunto de todos los objetos $x$ que
cumplen la propiedad $\varphi(x)$.

\begin{example}
    El conjunto de los números naturales impares se escribe
    \[
        \{x : x \in \omega \ \wedge \ x = 2k + 1 \ \text{para algún} \ k \in \omega\},
    \]
    y el conjunto de los números naturales pares menores que $10$ se escribe $\{x \in \omega : x < 10 \ \wedge \ x = 2k \
        \text{para algún} \ k \in \omega\} = \{0, 2, 4, 6, 8\}$.
\end{example}

\begin{remark}[El constructor sin restricciones]
    Empleado sin restricción alguna, el constructor $\{x : \varphi(x)\}$ conduce a contradicciones. Si se admite que
    \emph{toda} propiedad $\varphi$ determina un conjunto, nada impide formar el conjunto
    \[
        R = \{x : x \notin x\}
    \]
    de todos los conjuntos que no se pertenecen a sí mismos, y preguntar si $R \in R$. Ambas respuestas conducen a su
    contraria: si $R \in R$, entonces $R$ no cumple la propiedad que lo define y por tanto $R \notin R$; y si $R \notin R$,
    entonces $R$ sí cumple esa propiedad y por tanto $R \in R$. Esta es la \textbf{paradoja de Russell}\index{paradoja de
        Russell} (1901), y fue una de las razones que impulsaron la axiomatización de la teoría de conjuntos. En este capítulo
    usaremos el constructor de manera ingenua, como lo hizo Cantor; la restricción que lo vuelve seguro se estudiará más
    adelante.
\end{remark}

\section{Operaciones básicas sobre conjuntos}

Usando el constructor podemos definir algunas operaciones «algebraicas» básicas sobre dos conjuntos $X$ e $Y$:

\begin{itemize}
    \item la \textbf{intersección}\index{intersección} $X \cap Y = \{x : x \in X \ \wedge \ x \in Y\}$, cuyos elementos son los
          objetos que pertenecen a $X$ y a $Y$;
    \item la \textbf{unión}\index{unión} $X \cup Y = \{x : x \in X \ \vee \ x \in Y\}$, formada por los objetos que pertenecen a
          $X$ o a $Y$, o a ambos;
    \item la \textbf{diferencia}\index{diferencia} $X \setminus Y = \{x : x \in X \ \wedge \ x \notin Y\}$, formada por los
          elementos que pertenecen a $X$ pero no a $Y$;
    \item la \textbf{diferencia simétrica}\index{diferencia simétrica} $X \mathbin{\triangle} Y = (X \cup Y) \setminus (X \cap
              Y)$, formada por los elementos que pertenecen a la unión $X \cup Y$ pero no a la intersección $X \cap Y$.
\end{itemize}

En las fórmulas de la unión y la intersección usamos los conectivos lógicos $\wedge$ y $\vee$, que denotan las
operaciones lógicas \textbf{y} y \textbf{o}. El cuadro~\ref{tab:tabla-verdad} presenta la tabla de verdad de estas dos
operaciones y de otras cuatro: la \textbf{negación}\index{negación} $\neg$, la \textbf{implicación}\index{implicación}
$\Rightarrow$, la \textbf{equivalencia}\index{equivalencia} $\Leftrightarrow$ y la \textbf{barra de
    Sheffer}\index{barra de Sheffer} $\mid$, llamada también \textbf{nand}.

\begin{table}[htbp]
    \centering
    \caption{Tabla de verdad de los conectivos lógicos básicos.}
    \label{tab:tabla-verdad}
    \begin{tabular}{cc cccccc}
        \toprule
        $x$ & $y$ & $x \wedge y$ & $x \vee y$ & $x \Rightarrow y$ & $x \Leftrightarrow y$ & $x \mid y$ & $\neg x$ \\
        \midrule
        0   & 0   & 0            & 0          & 1                 & 1                     & 1          & 1        \\
        0   & 1   & 0            & 1          & 1                 & 0                     & 1          & 1        \\
        1   & 0   & 0            & 1          & 0                 & 0                     & 1          & 0        \\
        1   & 1   & 1            & 1          & 1                 & 1                     & 0          & 0        \\
        \bottomrule
    \end{tabular}
\end{table}

Tenemos entonces cuatro operaciones básicas sobre conjuntos $X$ e $Y$:
\[
    \begin{aligned}
        X \cap Y      & = \{x : x \in X \ \wedge \ x \in Y\},    & X \cup Y                           & = \{x : x \in X \ \vee \ x \in Y\}, \\
        X \setminus Y & = \{x : x \in X \ \wedge \ x \notin Y\},
                      & X \mathbin{\triangle} Y                  & = (X \cup Y) \setminus (X \cap Y).
    \end{aligned}
\]
